\section{Stato dell'arte e background}

\subsection{Bluesky e ATProto}

ATProto e' progettato come protocollo componibile per applicazioni sociali. La documentazione ufficiale descrive l'uso di lexicon comuni, repository personali, stream di eventi e servizi applicativi che possono essere implementati da attori diversi \parencite{atprotoOverview, blueskyAtprotoDocs}. A differenza di piattaforme centralizzate tradizionali, Bluesky separa piu' componenti dell'esperienza sociale: archiviazione, indicizzazione, visualizzazione, feed, labels e moderazione. \textcite{balduf2024looking} offrono la prima analisi larga scala della piattaforma, studiandone crescita, attivita' e diversita' dei provider, e rappresentano un riferimento diretto per la struttura dati usata nel presente lavoro.

\subsection{Moderazione componibile}

La documentazione Bluesky descrive la moderazione come un insieme di sistemi sovrapponibili: network takedowns, labels prodotte da servizi di moderazione, e controlli utente come mutes e blocks \parencite{blueskyLabelsModeration}. Questa separazione rende possibile una moderazione piu' pluralistica, nella quale il servizio ufficiale non e' l'unico attore a produrre segnali. La stessa architettura di labeling e' pensata per permettere a servizi indipendenti di pubblicare labels, che poi possono essere interpretate dai client o dagli AppView \parencite{blueskyModerationArchitecture}.

\subsection{Blocchi individuali}

Il blocco individuale e' un'azione pubblica: impedisce interazioni e nasconde l'account bloccato dall'esperienza dell'utente che blocca \parencite{blueskyBlocking}. Questa pubblicita' del record rende i blocchi osservabili come segnale comunitario. In letteratura, \textcite{bono2026selfmoderation} analizzano piu' di 100 milioni di azioni di blocco su Bluesky e mostrano che il blocking puo' essere studiato con profili comportamentali e modelli predittivi. Il presente lavoro usa i blocchi in modo diverso: non predice la propensione a essere bloccati, ma valuta se i blocchi ricevuti prima dell'evento contengano informazione sul successivo takedown ufficiale.

\subsection{Modlist e blocklist}

Le liste utente su Bluesky possono essere curate per feed o usate come modlist. La documentazione definisce \texttt{app.bsky.graph.defs\#modlist} come lista di utenti destinata a muting o blocking \parencite{blueskyUserLists}. Le modlist sono quindi una forma di moderazione comunitaria aggregata: un utente o una comunita' costruisce una lista e altri utenti possono adottarla. \textcite{jhaver2018blocklists} mostrano, nel caso di Twitter, che le blocklist possono proteggere da molestie ma introducono anche problemi di contestazione e accountability. Nel contesto Bluesky, \textcite{sokoto2026open} analizzano individual blocks e blocklists su larga scala, mostrando che le blocklist sono centrali ma eterogenee e concentrate.

\subsection{Labels e labeler indipendenti}

Le labels sono metadati pubblicati da servizi di moderazione; possono applicarsi ad account, record o altri oggetti, e sono interpretate secondo definizioni che determinano se contenuti o account siano nascosti, sfocati, avvertiti o ignorati \parencite{blueskyLabelsModeration}. I labeler indipendenti permettono una moderazione non esclusivamente centralizzata. Questa apertura e' rilevante per il presente lavoro perche' consente di confrontare labels ufficiali, labels di terzi e takedown ufficiali. Tuttavia, come si vedra' nei risultati, il problema empirico principale non e' solo il lift ma la coverage: un segnale molto concentrato nei positivi puo' restare poco utile se compare su una quota minima degli account.

\subsection{Account abusivi, spam e ban evasion}

La moderazione account-level e' spesso legata a problemi di autenticita', spam, impersonation e ban evasion. La letteratura su social bots e account falsi mostra che molte forme di abuso possono essere individuate tramite segnali di attivita', rete e metadati account-level \parencite{ferrara2016rise, varol2017online, cao2012aiding, stringhini2010detecting}. I transparency report di Bluesky collocano questo tema dentro la pratica della piattaforma: nel 2023 vengono menzionati strumenti proattivi contro slurs negli handle e spam accounts, mentre nel 2025 Bluesky descrive un approccio ibrido che combina automazione, sistemi anti-spam e revisione umana \parencite{bluesky2023moderation, bluesky2025transparency}. Queste fonti non identificano la causa dei singoli takedown analizzati qui, ma aiutano a interpretare con prudenza la dinamica del corner 00.

\subsection{Lavori empirici precedenti}

Il progetto si colloca tra tre filoni. Il primo e' lo studio delle architetture decentralizzate e dei protocolli sociali aperti, in cui Bluesky e ATProto sono casi recenti \parencite{balduf2024looking}. Il secondo e' lo studio della moderazione comunitaria e della self-moderation, incluso il ruolo di blocchi, blocklist e strumenti comunitari \parencite{seering2020reconsidering, jhaver2018blocklists, bono2026selfmoderation, sokoto2026open}. Il terzo e' lo studio di enforcement, sospensioni e deplatforming, dove la rimozione o limitazione account-level e' analizzata come decisione di piattaforma con effetti sociali e metodologici non riducibili al singolo contenuto \parencite{gillespie2018custodians, west2018censored, jhaver2021deplatforming}. Rispetto a questi lavori, il presente studio isola una domanda piu' specifica: quanto i segnali pubblici pre-takedown separano account successivamente rimossi da controlli negativi comparabili?
